\scanpage{307}
\mspage{(6)}
s'agit de désigner l'avatar ``objet simplicial'' ou ``complexe de chaines'' :
\[
 \begin{aligned}
 K_.&=\mathrm{ND}_ .(K_*), &K_*&=\mathrm{DP}_*(K_.)\\
 K_*&=(n\unsure{\mapsto|-}K_{[n]}), &K_.&=(K_n,d_n)_{n\geq0}
 \end{aligned}
\]
On est ainsi amené à noter par une notation $K_.$ \unsure{tt|tel} complexe de chaines, par $K_n$ ses composantes (ce qui exclut l'utilisation de la notation $K_n$ pour la $n$.\unsure{ème|ième} composante \inserted{de} $K_*$, et nous oblige à adopter une notation telle que $K_{[n]}$ !).

De $=$, le \unsure{topoïs|topos|topo|théorie} de DP appliqué à $A^o$ nous donne des équivalences
\diagram{\begin{tikzcd}[ampersand replacement=\&]
 A^* \arrow[r,shift left=1.2ex,"{\substack{\mathrm{ND}^{.}\\\approx}}"]
 \& A^{.}\arrow[l,shift left=1.2ex,"{\approx}"',"{\mathrm{DP}^{*}}"]
\end{tikzcd}}
et des notations
\[
 \begin{aligned}
 K^{.}&=\mathrm{ND}^{.}(K^*), & K^*&=\mathrm{DP}^*(K^{.})\\
 K_{\deleted{\unsure{*}}}^{*}&=(n\mapsto K^{[n]}), &K^{.}&=(K^n,d^n)_{n\geq0}.
 \end{aligned}
\]

Les foncteurs $\mathrm{ND}^{.}$ et $\mathrm{DP}^*$ commutent aux foncteurs additifs. En particulier, dans le cas d'une équivalence
\[
 A\underset{v}{\overset{\approx}{\longleftrightarrow}}B^o
\]
déjà \unsure{envisagé|envisagée}, étendant aux complexes de chaines ou de cochaines la \unsure{convent.|conv.|convention} de la montée ou de la descente des indices, on \unsure{trouve} que les foncteurs ND et DP y commutent. Donc dans ce cas, la $=$ lettre $K$ désigne
% archive sheet number: 308
